\chapter{Differentiation}
\section{Differentiation on $\mathbb{R}^n$}
\subsection{Continuity in MS}
\begin{remark}
    continuity 的定义: \textbf{open set 的 preimage 只能是 open set；等价于 closed set 的 preimage 只能是 closed set}
\end{remark}


\begin{theorem}
任意\textbf{ topological space 之间的连续映射} $ f:X \rightarrow Y$，\textbf{compact set 的 image 一定 compact}.
\end{theorem}

\begin{theorem}
任意 \textbf{topological space 之间的连续映}射 $f:X \rightarrow Y$，如果 \textbf{X compact，那么这个映射是 uniformly ctn 的}.（逆命题不成立）
\end{theorem}

\begin{definition}
    \textbf{$l_p$ norm on vector spaces:}\\
$$||v||_p = (\sum_{i = 1}^n |v_i|^p)^{1\over p}$$
\end{definition}

\begin{theorem}
    在\textbf{有限维 real vector space} （真包含 $\mathbb{R}^n$）中，\textbf{任意 $l_p$ norm 都等价}（无限维则不等).\\
    具体即：
    \begin{enumerate}
        \item 对于任意 $p_1, p_2$，都存在 $C_1, C_2$ 使得对于所有 $v\in V$:
        $$
        C_1||v||_{p_1} \leq ||v||_{p_2} \leq C_2||v||_{p_1}
        $$
        \item 任意 $l_p$ norm (通过差 norm 作为 metric) 定义出的\textbf{有限维 topological vector space 都是等价的}，于是它们在 continuity, differentiability 以及 convergence 等等概念上的形式都是一样的.
    \end{enumerate}
\end{theorem}

\subsection{Derivative}
\noindent 阐明了这件事情之后我们看欧式空间之间 differentiability 的定义：一个 $f: U \subset \mathbb{R}^n \rightarrow \mathbb{R}^m$ 在某个点 $x_0$ 上 diffble 的定义就是 \textbf{在 $x_0$ 上可以被某个 $Hom(\mathbb{R}^n, \mathbb{R}^m)$ 上的某个 Linear map 一阶近似，并且这个近似在 x 点附近 converge by norm（误差的下降速度 $>>$ 邻域的收缩速度）}
即：
\begin{definition}{Derivative} \label{derivartive}
    $$
    Df(x_0) \in \mathbb{R}^{m \times n} \text{ is the matrix } A \text{ such that } \lim_{||h|| \rightarrow 0} \frac{||f(x_0 + t) - f(x_0) - Ah||}{||h||} = 0
    $$
    注意这里的 $Df$ operator 是一个函数:
    $$
    Df : E \rightarrow Hom(R^n, R^m)
    $$
\end{definition}

我们可以由原函数和 derivative induce 出一个 remainder 函数，来表示在某点处 derivative 的 linear approximation 的误差程度随微小量 $h$ 变化的变化. 

\begin{definition}
我们定义 remainder function: (用来做分析时使用)
\begin{align}
    r_{x_0} : \mathbb{R}^n &  \rightarrow \mathbb{R}^m \\
    h & \mapsto f(x_0 + h) - f(x_0) - Df(x_0) \cdot h
\end{align}
\end{definition}

\noindent 对于任意的 $x_0$，总是有 $||\frac{r_{x_0}(h)}{h}|| \rightarrow 0$ as $h \rightarrow 0$，因而 $r(h)$ 是 $o(h)$ 的.

\begin{theorem}{$\mathbb{R}^n \rightarrow \mathbb{R}^m$ is similar to $\mathbb{R}^n \rightarrow \mathbb{R}$} \label{$f$ ctn/diff iff $f_i$ ctn/diffble for all $i$}
    For $$f: \mathbb{R}^n \rightarrow \mathbb{R}^m = \begin{pmatrix}
        f_1 \\ f_2 \\ \vdots \\ f_m
    \end{pmatrix}
    $$ we always have
    $f$ ctn/diff iff $f_i$ ctn/diffble for all $i$
\end{theorem}
意思是在证明可导/连续时我们总是可以只考虑 $m=1$ 的情况. 这件事非常有用.

\section{Directional Derivative}
\begin{definition}
    固定一个 $u \in \mathbb{R}^n$，define 
    \begin{align}
        D_u f(X_0) : \mathbb{R}^n & \rightarrow \mathbb{R}^m\\
        x & \mapsto \lim_{t\rightarrow 0, t\in \mathbb{R}} \frac{f(x + tu) - f(x)}{t}
    \end{align}

        Notice that:
        $$
        D_u f(x_0) = \frac{d}{dt}\bigg| _{t=0} f(x_0 + tu)
        $$
\end{definition}


\begin{proposition}{Intuition of Directional Derivative}
我们发现，实际上 $D_v f$ 和 $f$ 是来自同一个 function space $\mathcal{F}(\mathbb{R}^n, \mathbb{R}^m)$ 的.\\
实际上，$D_v f$ 描述的对于每个 $x \in dom(f)$，在点 $x$ 处 $f$ 在 $v$ 这个方向上的变化率. 这里 $f$ 才是变量，$v$ 是固定的。\\
$D_v f$ 对每个(存在该 directional derivative)的 $x \in dom(f)$ 都给出了一个 $\mathbb{R}^m$ 数值，对应 $f(x)$ 的值给出了 $f$ 在方向 $v$ 上的变化率 $D_v f(x)$. 所以整个 $D_v f$ 就是 $f$ 对应的一个变化率函数.
\end{proposition}

\begin{remark}
    $\mathbb{R}^n$ 是一个 vector space. 我们自然想到，如果知道每个 $e_i$ 方向上 $f$ 对应的变化率函数 $D_{e_i} f$，那么就知道了整体 $f$ 在各个方向上的变化率函数. 我们对这个函数投一个点 $x = c_1 e_1 + \cdots + c_n e_n \in dom(f)$:
    $$
    Df(x_0)(x) = \sum_{i=1}^n c_i D_{e_i} f(x_0) 
    $$
    就可以得到这点上 $Df(x_0)$ 的作用.
\end{remark}

\noindent 我们有一个专门的称呼来指代 $D_{e_i} f$
\begin{definition}{partial derivative} \label{partial derivative}
    我们使用 $\frac{\partial }{ \partial x_i} f$ 来表示 $D_{e_i} f$
\end{definition}

接着刚才的 remark，我们发现
\begin{theorem}
    如果 $f $ 在 $x_0 \in dom(f) $ 处 diffble，那么它在 $x_0$ 处的 derivative 为：
    \begin{align}
        Df(x_0) 
        & = \begin{pmatrix}
        \frac{\partial f_1}{\partial x_1} (x_0)& \frac{\partial f_1}{\partial x_2}(x_0) & \cdots & \frac{\partial f_1}{\partial x_n}(x_0) \\
        \frac{\partial f_2}{\partial x_1} (x_0)& \frac{\partial f_2}{\partial x_2}(x_0) & \cdots & \frac{\partial f_2}{\partial x_n}(x_0) \\
        \vdots & \vdots & \ddots & \vdots \\
        \frac{\partial f_m}{\partial x_1}(x_0) & \frac{\partial f_m}{\partial x_2} (x_0)& \cdots & \frac{\partial f_m}{\partial x_n}(x_0)
        \end{pmatrix}
    \end{align}

    其中每一列就是 $D_{e_i} f(x_0) = \frac{\partial}{\partial x_i} f(x_0)$. 因而对于每个输入 $ c \in \mathbb{R}^n$，我们得到的输出都是 $\sum_{i=1}^n c_i \frac{\partial}{\partial x_i} f(x_0) \in \mathbb{R}^m$.\\
    
    \begin{align}
        Df(x_0) (c)
        & = \begin{pmatrix}
        \sum_{i=1}^n c_i \frac{\partial}{\partial x_i} f_1(c) \\
        \sum_{i=1}^n c_i \frac{\partial}{\partial x_i} f_2(c) \\
        \vdots  \\
        \sum_{i=1}^n c_i \frac{\partial}{\partial x_i} f_n(c)
        \end{pmatrix}
    \end{align}
我们称这个 matrix 为 $f$ 的 \textbf{Jacobian matrix at $x_0$}
\end{theorem}

\begin{remark}
    $$
        Df: A \subset \mathbb{R}^n \rightarrow \text{Hom}(\mathbb{R}^n, \mathbb{R}^m) 
    $$
    $$
        x \mapsto 
        \begin{pmatrix}
        \frac{\partial f_1}{\partial x_1} (x)& \frac{\partial f_1}{\partial x_2}(x) & \cdots & \frac{\partial f_1}{\partial x_n}(x) \\
        \frac{\partial f_2}{\partial x_1} (x)& \frac{\partial f_2}{\partial x_2}(x) & \cdots & \frac{\partial f_2}{\partial x_n}(x) \\
        \vdots & \vdots & \ddots & \vdots \\
        \frac{\partial f_m}{\partial x_1}(x) & \frac{\partial f_m}{\partial x_2} (x)& \cdots & \frac{\partial f_m}{\partial x_n}(x)
        \end{pmatrix}
    $$
    \textbf{显然，如果在 $A$ 上每个 partial derivative 都是 continuous 的，那么这个 $Df$ 也是 continuous 的（if we define norm on $\text{Hom}(\mathbb{R}^n, \mathbb{R}^m) $}
\end{remark}

\begin{definition}{higher-order partial derivatives}
    我们知道 $f: \mathbb{R}^n \rightarrow \mathbb{R}^m$ 的任意 partial derivative ${\partial \over \partial x_i} f: \mathbb{R}^n \rightarrow \mathbb{R}^m$ 仍然是一个 $\mathbb{R}^n \rightarrow \mathbb{R}^m$ 的函数. 我们再对这个函数求 $e_j$ 方向上的 partial derivatve，可以得到另一个 ${\partial^2 \over \partial x_j \partial x_i} f: \mathbb{R}^n \rightarrow \mathbb{R}^m$.\\
    recursively define: 
    $$
    {\partial^d \over \partial x_{k_d}\partial x_{k_{d-1}}\cdots \partial x_{k_1}} = \frac{\partial }{\partial x_{k_d}}[{\partial^{d-1} \over \partial x_{k_{d-1}}\partial x_{k_{d-2}}\cdots \partial x_{k_1}}]
    $$
    为一个 $d^\text{th}$-oder 的 partial derivative of $f$
\end{definition}
\begin{remark}
    只有 partial derivative 是有多阶的，而 derivative $D$ 也急速hi全导数是没有多阶的，因为它本身就是一个 linear map. 因为 $f: \mathbb{R}^n \rightarrow \mathbb{R}^m$ 的任意阶段偏导数也是 $\mathbb{R}^n \rightarrow \mathbb{R}^m$ 的，但是全导数是一个从元素到线性变换的函数.\\
    我们平时使用的单变量real-valued函数的高阶导数是 partial derivative 的特殊情况，这是在 $\mathbb{R} \rightarrow \mathbb{R}$ 上的函数只有一个变量，全导数和偏导数是同一个函数.
\end{remark}

\section{$C^r$ class}
\begin{definition}{$C^r$ class}
    如果在 $A \subset \text{dom}(f)$ 上， $f$ 的所有 $r^\text{th}$-order partial derivatives 都 \textbf{exist 并且 ctn}，则称 $f \in C^r(A)$.
    \\如果对于任意 $r \in \mathbb{N}$ 都有 $f \in C^r(A)$，则称 $f \in C^{\infty}$
\end{definition}


\begin{remark}
    $f \in C^{r+1}(A) \Leftrightarrow $ all partials of $f$ are in $C^r(A)$
\end{remark}



    如果 f 在点 $x$ 附近 being $C^1$: $f$ 在点 $x$ 附近，每个 $e_i$ 切向量上的变化率函数都是连续的. 所以我们有一种直觉：$f$ 在这个点上是可以被一个 linear map 近似的.

\begin{theorem}{Sufficient condition for $f$ differentiable at some point}
    如果 $f: \mathbb{R}^n \rightarrow \mathbb{R}^m$ 在点 $x_0$ 附近是 $C^1$ 的，那么它在 $x_0$ 上是 differentiable 的 (即 $Df(x_0)$ exists.\\
    (Note: 此时在 $x_0$ 附近 $Df: A \subset \mathbb{R}^n \rightarrow \text{Hom}(\mathbb{R}^n, \mathbb{R}^m) $ 也是 continuous 的.)
\end{theorem}
\begin{proof}
见L05.
\end{proof}

\begin{remark}
    因而如果 $f \in C^1(dom(f))$, $f$ 就是 differentiable 的函数.
\end{remark}


\begin{theorem}{Mixed partials}
    If $f$ is $C^2$ 则任取 $x \in \text{dom}(f)$，$f$ 在 $x$ 处的任意 $2^\text{nd}$-order partial 可以更换次序. 即:
    $$
    {\partial^2 \over \partial x_{j_1} \partial x_{j_2}}f(x_0) = {\partial^2 \over \partial x_{j_2} \partial x_{j_1}} f(x)
    $$
\end{theorem}
\begin{proof}
    见 L06.\\
    对于这种 multivariable 的关于 diffbility/ ctnity 的证明，我们总是应该想到把它先等价为 $\mathbb{R}^n \rightarrow \mathbb{R}$，然后再对 $n$ 个变量进行逐个处理，运用 single variable analysis 的各种 mean value theorem 来处理.
\end{proof}


\begin{corollary}{Mixed partials}\label{Mixed partials}
    自然可以得到：如果 $f \in C^r(A)$，那么 $f$ 的任意 $r$ 阶及以下阶 partial derivative 都可以换序. 即: $\forall m \leq r$ 以及 $\forall \pi \in S_m$:
    $$
     {\partial^m \over \partial x_{k_1}\partial x_{k_{2}}\cdots \partial x_{k_m}} = {\partial^m \over \partial x_{k_{\pi(1)}}\partial x_{k_{\pi(2)}}\cdots \partial x_{k_{\pi(m)}}}
    $$
    因而如果 $f\in C^\infty (A)$ 那么所有的 partial derivative 都可以换序.
\end{corollary}


\section{Chain Rule, Product Rule and Taylor's theorem}
\begin{definition}{multi-index notation} \label{multi-index}
一个 multi-index 是一个 $n-$tuple，其中每个 entry 都是一个 non-neg integer.\\
我们使用以下四个 notation 来进行简写:\\
$$
|\alpha| = \sum_{j = 1}^n \alpha_j
$$
$$
\alpha! = \prod_{i = 1}^n (\alpha_j !)
$$
$$
x^{\alpha} = {x_1}^{\alpha_1}{x_2}^{\alpha_2} \cdots {x_n}^{\alpha_n}
$$
$$
\partial ^ \alpha f ={\partial^{|\alpha|} \over (\partial {x_1})^{\alpha_1} (\partial {x_2})^{\alpha_2} \cdots (\partial {x_n})^{\alpha^n}} f  
$$
\end{definition}
这个 notation 的引入会使得很多定理的表述简洁很多.
\begin{example}
$$
\partial ^{(2,1)} f = \frac{\partial^3}{\partial x_1 \partial x_2 \partial x_3 } f
$$
$$
(3,3,0)! = 3!3!0! = 36
$$
\end{example}


\subsection{Chain Rule}
\noindent Recall chain rule in single variable analysis:
$$f: A \subset \mathbb{R} \rightarrow B \subset \mathbb{R} \;\; \text{, } g: B \rightarrow \mathbb{R} 
$$
with A, B open; $f$ is diffble at some point $x$ and $g$ is diffble at $f(x)$.\\
Then we have: 
$$
\frac{d}{dx} (g \circ f) (x) = g'(f(x)) f'(x)
$$

现在我们有 multi 的版本:
\begin{theorem}
    如果 $$f: A \subset \mathbb{R}^n \rightarrow B \subset \mathbb{R}^m \;\; \text{, } g: B \rightarrow \mathbb{R}^p 
$$
with A, B open; $f$ is diffble at some point $x$ and $g$ is diffble at $f(x)$.\\
Then we have:
$$
D (g \circ f) (x) =Dg(f(x))\, Df(x)
$$
\end{theorem}
和 single varable 的表述基本一模一样.
\begin{proof}
    见 L07. 证明思路即取 $f,g, g\circ f$ 的 remainder function 并发现在 $R_f, R_g$ 趋向 0 时 $R_{g\circ f}$ 也趋向 0.
\end{proof}

\subsection{Multinomial Theorem}
\begin{theorem}{Multinomial Thm} \label{Multinomial Thm}
$\forall (x_1, x_2, \cdots, x_n) \in \mathbb{R}^n$ 以及任意 $k \in \mathbb{Z}_{\geq 0}$ we have:
$$
{(x_1 + \cdots + x_n)}^k  = \sum_{|\alpha| = k} \frac{k!}{\alpha !} x^\alpha 
$$
\end{theorem}
这是一个 combinatorics 的 proposition. 
\begin{proof}
    (Informal proof, but sufficient to show it.)\\
    ${(x_1 + \cdots + x_n)}^k$ 每项就是每次在 $\{x_1, x_2, \cdots, x_n\}$ 中取一个元素，取 $k$ 次得到，一共有 $\sum_{|\alpha| = k} \frac{k!}{\alpha !} = k!$ 个项，每一项都是一个 $x^\alpha$ where  $|\alpha| = k$.\\
    每个 $x^\alpha$ 会被取到重复次，这个重复的数量为 ${k! \over \alpha !}$\\
     因为 $k! \over \alpha !$ 就是把一个 size 为 $k$ 的 set 分成 size 分别是 $\alpha_1, \cdots, \alpha_n$ 的子集的方法.\\
     推导: 
$$
\# =  C^k_{\alpha_1} \times  C^{k-\alpha_1}_{\alpha_2}  \times C^{k-\alpha_1 - \alpha_2}_{\alpha_3} \times\cdots \times C^{k - \sum_{i=1}^{n-1} \alpha _i}_{\alpha_n} = {k! \over \alpha_1! \alpha_2 ! \cdots \alpha_n !}
$$
    （对于 for example:  $\alpha = \{1,0,2,2,0\}$ 即 $x_1 {x_3}^2 {x_4}^2$ 这一项，会被取到 ${k! \over \alpha !} = \frac{5!}{2!2!}$ 次. 这个次数等于$ 5 \times C^4_2 \times C^2_2 $ ）
    \\ formal 的证明即 by induction，使用 binomial thm 即可.
\end{proof}

\subsection{Product Rule}

\begin{lemma}{在 $\mathbb{R} \rightarrow \mathbb{R}$ 上的函数的 product rule}
令 $f_1, \cdots f_m: \mathbb{R} \rightarrow \mathbb{R}$

$$
\frac{d^k {(f_1,\cdots f_m)}}{dx^k} = \sum_{|\alpha| = k} \frac{k!}{\alpha !} (\frac{d^{\alpha_1} f_1}{dx^{\alpha_1}}) \cdots (\frac{d^{\alpha_m} f_m}{dx^{\alpha_m}})
$$

\end{lemma}

\begin{note}
    特别地，$$
    \frac{d^k}{dx^k} (fg) = \sum_{a+b = k} \frac{k!}{a!b!} (\frac{d^a}{dx^a}f) (\frac{d^b}{dx^b} g)
    $$
    这个式子也等价于 $$
    \sum_{a = 0}^k \frac{k!}{a!(k-a)!} (\frac{d^a}{dx^a}f) (\frac{d^{k-a}}{dx^{k-a}} g)
    $$
\end{note}

\begin{proof}
    见 HW5 G. \\
\end{proof}

\begin{theorem}{product rule} \label{product rule}
$\forall \alpha \in \mathbb{Z}_{\geq 0}^n$，$f,g: \mathbb{R}^n \rightarrow \mathbb{R}$, 如果 $f,g \in C^{|\alpha|} (\mathbb{R}^n)$, then：
$$
\partial^{\alpha} (fg) = \sum_{\beta + \gamma = \alpha} \frac{\alpha !}{\beta ! \gamma !} (\partial ^ \beta f) (\partial ^ \gamma g)
$$
\end{theorem}
\begin{proof}
    由 lemma 由归纳得到，见 L08.
\end{proof}
\begin{note}
    实际上这里 f,g 是 $\mathbb{R}^n \rightarrow \mathbb{R}$ 的函数并不是重要的，重要的是 $\alpha$ 和几个变量有关系。如果 $f,g:\mathbb{R}^{100} \rightarrow \mathbb{R} $，$\alpha \in \mathbb{Z}_{\geq 0} ^{100}$ 但是只有一个 entry 不是 0，那么我们可以看作 $\alpha \in \mathbb{Z}_{\geq 0}^1$，使用 $\mathbb{Z}_{\geq 0}^1$ 上的 Product rule.
\end{note}


\begin{lemma}{Taylor Thm in $\mathbb{R}$}
如果 $f: I \subseteq \mathbb{R} \rightarrow \mathbb{R} \in C^k(I)$，并且 f 是 $(k+1)$ - diffble 的.\\
(即，$f$ 几乎是 $C^{k+1}$ 的，但是并不需要 k+1 阶导数也连续，条件更宽)\\
任取 $a \in I$，我们定义 \textbf{the $k^\text{th}$ Taylor polynomial centered at $a$}:
$$
T_{a,k}(x) = f(a) + \sum_{j=1}^k \frac{f^{(j)}(a)}{j!} (x-a)^j
$$
Then: 这个 $T_{a,k}(x)$ 是 \textbf{best $k^\text{th}$-order polynomial approximation centered at $a$}: $$
f(x) - T_{a,k}(x) =  R_{a,k} (x) = \frac{f^{(k+1)}(c_x)}{(k+1)!} (x-a)^{k+1} = o(|x-a|^k)
$$
Note: 这里的 remainder 和 a,k 都有关系. 是：
对于每个 $x \in I$，$c_x \in [a,x]$ 是由 Cauchy MVT 决定的一个中间值.
\end{lemma}

\begin{remark}
   但是这仍然并不能说明如果 $f \in C^{\infty}(I)$，就一定有 $f(x) = T_{a, \infty}(x), \forall x,a \in I$. \\
   Note:  $C^\omega (\mathbb{R}) \subset C^\infty(\mathbb{R}).$\\
    一个光滑函数 $f$ 在 $I$ 上解析的条件是在 $I$ 上这个余项函数是处处收敛到 0 的.
    具体: $
    \forall \text{compact } K \subseteq I, \; \exists m \in \mathbb{R} \text { s.t. } $
    $$
    \bigg |\frac{d^j}{dx^j} f(x) \bigg | \leq m^{k+1} j! 
    $$
    for all $x \in K$ and for all order $j$.
\end{remark}

\begin{proof}
    具体见 L08.\\
    思路：Recall that Rolle MVT $\rightarrow$ Lagrange MVT $\rightarrow$ Cauchy's MVT\\
    我们使用 Cauchy's MVT apply to 余项来证明 Taylor's Thm.
\end{proof}


\begin{theorem}{Taylor's Thm}\label{Taylor's Thm}
    
\end{theorem}











\section{Inverse function Theorem}
\begin{theorem}{Inverse function Theorem} \label{IFT}
    如果 $f: A \subseteq \mathbb{R}^n \rightarrow \mathbb{R}^m$ 是 $C^r$ 
 ($r \geq 1$)的，对于任意 $x_0 \in A$ s.t. $\det{Df(x_0)} \not= 0$，我们都有：
    \begin{enumerate}
        \item 存在某个 $x_0$ 的 nbh $U$ 以及 $f(x_0)$ 的 nbh $V$ 使得 $f(U) = V$ 且 $f|_U : U \rightarrow V$ 是 bijective 的. (即 $f$ 在 $x_0$ 附近是 \textbf{locally bijective }的).
        \item $f|_U$ 的 inverse function $g: V \rightarrow U$ 也是 $C^r$ 的，并且对于 $U$ 中任意 $x$，都有 $Dg(f(x)) = (Df(x))^{-1}$ \\ 
        (这是显然的，因为 $I_n = D (g \circ f) (x) = Dg(f(x))Df(x)$)
    \end{enumerate}
\end{theorem}
\begin{remark}
    比较直观。我们知道如果 $f \in C^1(A)$，意味着 Jacobian matrix 中每一项都是连续的函数，那么\textbf{整个 $Df: \mathbb{R}^n \rightarrow \mathbb{R}^{n \times n}$ 在 $x_0$ 附近也是连续的}；并且由于 \textbf{$\det :\mathbb{R}^{n \times n} \rightarrow \mathbb{R}$ determinant 函数是连续的}，我们知道整个 $Df$ 在 $x_0$ \textbf{locally 都是 non-sigular 的}. \\\\
    我们可以类比 single variable 情况下：如果 $C_1$ 的函数 $f$ 的导函数 $f'$ 在某处不是 0，那么在这一块地方附近 $f$ 都 monotonic. 这是因为 $f'$ 连续，所以这附近 $f'$ 都非 0 且同向.\\\\
    IFT 告诉我们的是：这个性质可以推广到多元函数，by: 单元函数导数值不为 0 扩展到多元函数导数的 det 不为0.\\
    并且 IFT 还多告诉我们的一件事情是: 这个 local inverse function 会保留 $C^r$ 的性质.\\
\end{remark}

\begin{lemma}
    对于任意的 matrix $C \in \mathbb{R}^{m \times n}$，并任取 $v \in \mathbb{R}^n$，我们都有:
    $$
    |Cx| \leq ||C|| \cdot |x|
    $$
    并且： 
    $$
||C||\leq nm \max_{1 \leq i \leq m, 1\leq j \leq n} |C_{ij}|
$$
where we take the sup norm of a linear map:
$$
||C|| = \sup_{x \in \mathbb{R}^n, |x| = 1} |Cx|
$$
\end{lemma}


\begin{lemma}
    如果 $f: A \rightarrow \mathbb{R}^n$ 是 $C^1$ 的，且 $Df(a)$ non-singular，即 $\det(Df(a)) \not = 0$，那么存在某个 $a$ 的 nbh $B_{\delta} (a)$ 使得对于任意 $x,y \in B_{\delta}(a)$，都有：
    $$
    |f(x_0) - f(x_1)| \geq \alpha |x_0 - x_1|
    $$
\end{lemma}
\begin{note}
 \textbf{这是一个比 locally injective 更强的条件.}
\end{note}







\section{Implicit Function Thm}












